%% ========================================================================
%% Chapter 14: Backup dan Pemulihan Sistem
%% ========================================================================

\chapter{Backup dan Pemulihan Sistem}
\label{ch:backup-and-recovery}

\chapterhero{gray}{Rancang cadangan data yang andal, otomasi jadwal backup, dan latihan pemulihan agar kehilangan data tidak berubah menjadi bencana.}{rsync}{cron}{pemulihan}

Bayangkan kamu bekerja seharian penuh menyelesaikan laporan penting, kemudian
laptop mati mendadak dan semua file hilang. Skenario ini bukan sekadar cerita
menakutkan: kegagalan disk keras, kesalahan penghapusan file, serangan malware,
atau sekadar mati listrik di saat yang salah adalah kejadian nyata yang dialami
oleh administrator sistem di seluruh dunia setiap harinya. Backup adalah
jaring pengaman yang memastikan kejadian buruk tidak berubah menjadi bencana.

Di lingkungan Linux, ada tiga alat utama yang menjadi tulang punggung sistem
backup. \cmd{rsync} bekerja seperti sinkronisasi dokumen antara dua folder:
efisien, cerdas, dan bisa berjalan melalui jaringan. \cmd{tar} bekerja seperti
memadatkan dokumen ke dalam amplop arsip yang bisa disimpan di lemari: semua
file dibungkus rapi menjadi satu arsip terkompresi. \cmd{dd} bekerja di level
yang lebih mendasar: menyalin setiap bit dari satu disk ke disk lain, seperti
membuat kloning kartu identitas beserta semua isinya secara bit per bit.

Memilih alat yang tepat bergantung pada apa yang perlu di-backup, seberapa
sering, dan seberapa cepat pemulihan harus bisa dilakukan. Bab ini membahas
ketiga alat tersebut, cara menjadwalkan backup secara otomatis, dan prosedur
pemulihan ketika skenario terburuk benar-benar terjadi.

\textbf{Yang Akan Kamu Pelajari}

Setelah menyelesaikan bab ini, kamu akan mampu:

\begin{enumerate}
    \item menjelaskan perbedaan backup full, incremental, dan differential
        serta menentukan strategi yang tepat untuk situasi tertentu;
    \item menggunakan \cmd{rsync} untuk sinkronisasi direktori lokal dan membuat
        backup snapshot berbasis hard link;
    \item membuat arsip backup dengan \cmd{tar} beserta opsi kompresi dan
        incremental yang sesuai;
    \item menjadwalkan skrip backup otomatis menggunakan \cmd{cron} dengan
        sintaks crontab yang benar;
    \item memulihkan file dan direktori dari backup serta memverifikasi
        integritasnya menggunakan \cmd{md5sum} dan \cmd{sha256sum}.
\end{enumerate}

\begin{notebox}
Seluruh praktek pada bab ini dijalankan pada direktori \dir{~/lab-os/chapter12-backup/}.
\begin{lstlisting}[language=bash]
mkdir -p ~/lab-os/chapter12-backup
cd ~/lab-os/chapter12-backup
\end{lstlisting}
\end{notebox}

%% ========================================================================
\section{Strategi Backup}
\label{sec:backup-strategy}
%% ========================================================================

Sebelum memilih alat, kamu perlu memiliki strategi. Strategi backup yang baik
menjawab tiga pertanyaan: apa yang perlu di-backup, seberapa sering, dan
disimpan di mana. Tanpa jawaban yang jelas untuk ketiga pertanyaan ini, backup
yang kamu lakukan mungkin tidak berguna saat dibutuhkan.

Ada tiga jenis backup yang umum digunakan, masing-masing dengan karakteristik
yang berbeda dalam hal ruang penyimpanan dan kecepatan pemulihan.

\emph{Full backup} menyalin seluruh data yang dipilih tanpa pengecualian.
Pemulihan paling sederhana karena hanya membutuhkan satu set backup. Kelemahannya
adalah ukuran besar dan waktu proses lama, sehingga tidak praktis dilakukan
setiap hari untuk data bervolume besar.

\emph{Incremental backup} menyalin hanya file yang berubah sejak backup
terakhir, baik full maupun incremental sebelumnya. Ukurannya kecil dan cepat,
tetapi pemulihan memerlukan rekonstruksi dari backup full ditambah semua
incremental secara berurutan.

\emph{Differential backup} menyalin file yang berubah sejak full backup
terakhir. Ukurannya lebih besar dari incremental tetapi pemulihan lebih
sederhana: hanya perlu backup full dan differential terbaru.

Strategi 3-2-1 adalah standar industri yang direkomendasikan untuk memastikan
backup tidak menjadi \emph{single point of failure}. Angka-angkanya mudah
diingat: simpan \textbf{3} salinan data, di \textbf{2} media berbeda, dengan
\textbf{1} salinan di lokasi fisik yang berbeda (offsite). Jika server kantor
terbakar, salinan offsite di cloud atau lokasi lain tetap aman.

Dua metrik kritis perlu ditentukan sebelum merancang sistem backup. \emph{RPO}
(Recovery Point Objective) adalah pertanyaan: seberapa banyak data yang boleh
hilang? Jika RPO adalah 24 jam, backup harus berjalan setidaknya sekali sehari.
\emph{RTO} (Recovery Time Objective) adalah pertanyaan: berapa lama downtime
yang bisa ditoleransi? Jika RTO adalah 2 jam, proses restore harus bisa selesai
dalam 2 jam sejak kegagalan terjadi.

\subsection*{Praktek 12.1: Rencanakan Strategi Backup}

Tujuan praktek ini adalah menyusun rencana backup tertulis menggunakan heredoc
dari Bab~3, lalu membuat data simulasi yang akan digunakan di praktek-praktek
berikutnya.

\begin{enumerate}
    \item Buat direktori struktur data simulasi yang akan di-backup:
\begin{lstlisting}[language=bash, caption={Menyiapkan direktori dan data simulasi}]
cd ~/lab-os/chapter12-backup
mkdir -p data-sumber/{dokumen,konfigurasi,log}

echo "laporan-keuangan-2026.txt" > data-sumber/dokumen/laporan.txt
echo "catatan-penting.txt" >> data-sumber/dokumen/laporan.txt
echo "app_port=8080" > data-sumber/konfigurasi/app.conf
echo "db_host=localhost" >> data-sumber/konfigurasi/app.conf
for i in $(seq 1 20); do
    echo "$(date) - log entry $i" >> data-sumber/log/app.log
done

du -sh data-sumber/*/
find data-sumber -type f | wc -l
\end{lstlisting}
    \textit{Amati:} berapa total file yang ada di direktori simulasi ini?

    \item Buat dokumen rencana backup menggunakan heredoc:
\begin{longlisting}[language=bash, caption={Membuat rencana backup dengan heredoc}]
cat > rencana-backup.txt << 'EOF'
====================================================
RENCANA BACKUP SISTEM LAB LINUX
====================================================
Tanggal Dibuat : (isi tanggal kamu)
Nama Sistem    : Lab Ubuntu Server 22.04
====================================================

INVENTARIS DATA
---------------
data-sumber/dokumen/    : laporan dan catatan, KRITIS
data-sumber/konfigurasi/: konfigurasi aplikasi, KRITIS
data-sumber/log/        : log sistem, TIDAK KRITIS

TARGET METRIK
-------------
RPO : 24 jam (data boleh hilang maksimal 1 hari)
RTO : 2 jam  (sistem harus pulih dalam 2 jam)

STRATEGI
--------
Full backup    : setiap Minggu malam (dokumen + konfigurasi)
Incremental    : Senin-Sabtu malam (dokumen + konfigurasi)
Log            : TIDAK dibackup (dapat diregenerasi)
Penyimpanan    : lokal ~/lab-os/chapter12-backup/arsip/

ATURAN 3-2-1
-----------
Salinan 1 : direktori data-sumber (data asli)
Salinan 2 : direktori arsip/ lokal (backup lokal)
Salinan 3 : server remote / cloud storage (offsite)
====================================================
EOF

cat rencana-backup.txt
\end{longlisting}

    \item Hitung estimasi kebutuhan ruang backup untuk 30 hari:
\begin{lstlisting}[language=bash, caption={Estimasi kebutuhan ruang backup}]
DATA_SIZE=$(du -sb data-sumber/dokumen/ data-sumber/konfigurasi/ \
    | awk '{sum+=$1} END{print sum}')
echo "Ukuran data kritis: $DATA_SIZE byte"
echo "Full backup x4/bulan: $((DATA_SIZE * 4)) byte"
echo "Incremental est. 10% x24/bulan: $((DATA_SIZE * 24 / 10)) byte"
echo "Total estimasi 30 hari: $((DATA_SIZE * 4 + DATA_SIZE * 24 / 10)) byte"
\end{lstlisting}
    \textit{Amati:} berapa perbandingan kebutuhan ruang strategi full+incremental
    dibandingkan jika full backup dilakukan setiap hari (full x30)?
\end{enumerate}

\begin{challengebox}
Buat file teks berisi rencana backup yang lebih lengkap menggunakan teknik
heredoc dari Bab~3. Rencana harus mencantumkan: jadwal harian dan mingguan,
estimasi ruang untuk setiap jenis backup, lokasi penyimpanan, dan prosedur
pengujian restore. Simpan ke file \file{rencana-lengkap.txt}. Kemudian gunakan
teknik redirection dan pipeline dari Bab~3 untuk menghitung berapa baris
rencana yang kamu tulis.
\end{challengebox}

%% ========================================================================
\section{rsync untuk Sinkronisasi dan Backup}
\label{sec:rsync-backup}
%% ========================================================================

\cmd{rsync} adalah alat sinkronisasi yang sangat efisien. Rahasianya terletak
pada algoritma delta-transfer: saat menyinkronkan file yang sudah ada di tujuan,
\cmd{rsync} hanya mengirimkan bagian yang berubah, bukan seluruh file dari awal.
Bayangkan mengedit dokumen Word: daripada mengirim ulang seluruh dokumen 10 MB,
\cmd{rsync} hanya mengirim paragraf yang kamu ubah.

Opsi \texttt{-a} (archive) adalah kombinasi dari beberapa opsi sekaligus:
\texttt{-r} (rekursif), \texttt{-l} (simpan symlink), \texttt{-p} (simpan
permission), \texttt{-t} (simpan timestamp), \texttt{-g} (simpan group), dan
\texttt{-o} (simpan owner). Hampir selalu digunakan bersama \texttt{-v} (verbose)
untuk melihat kemajuan.

\begin{lstlisting}[language=bash, caption={Penggunaan dasar rsync}]
rsync -av ~/lab-os/chapter12-backup/data-sumber/ ~/lab-os/chapter12-backup/arsip-rsync/
rsync -avn ~/lab-os/chapter12-backup/data-sumber/ ~/lab-os/chapter12-backup/arsip-rsync/
rsync -av --delete ~/lab-os/chapter12-backup/data-sumber/ ~/lab-os/chapter12-backup/arsip-rsync/
rsync -avz ~/data/ user@server-remote:/backup/data/
\end{lstlisting}

Opsi \texttt{-n} (dry run) mensimulasikan proses tanpa benar-benar mengubah
apa pun. Selalu jalankan dry run sebelum menggunakan \texttt{--delete}, karena
opsi itu akan menghapus file di tujuan yang tidak ada di sumber.

\begin{warningbox}
Perhatikan perbedaan trailing slash pada path sumber. \cmd{rsync -av
/home/user/} (dengan slash) menyalin isi direktori. \cmd{rsync -av /home/user}
(tanpa slash) menyalin direktori itu sendiri, menghasilkan subdirektori
tambahan di tujuan. Kesalahan trailing slash adalah sumber kebingungan
yang sangat umum.
\end{warningbox}

Fitur paling canggih \cmd{rsync} untuk backup adalah opsi \texttt{--link-dest}.
Dengan opsi ini, setiap snapshot harian terlihat seperti full backup lengkap
saat ditelusuri, tetapi file yang tidak berubah hanya berupa hard link ke file
snapshot sebelumnya. Tidak ada duplikasi data di disk.

\begin{lstlisting}[language=bash, caption={Backup snapshot berbasis hard link}]
BACKUP_BASE=~/lab-os/chapter12-backup/snapshots
TODAY=$(date +%F)
YESTERDAY=$(date -d yesterday +%F)

mkdir -p "$BACKUP_BASE/$TODAY"
rsync -av --delete \
    --link-dest="$BACKUP_BASE/$YESTERDAY" \
    ~/lab-os/chapter12-backup/data-sumber/ \
    "$BACKUP_BASE/$TODAY/"
\end{lstlisting}

Opsi \texttt{--exclude} dan \texttt{--exclude-from} mengecualikan file atau
pola tertentu dari sinkronisasi. Opsi \texttt{--log-file} mencatat semua
aktivitas ke file untuk audit.

\begin{lstlisting}[language=bash, caption={Opsi rsync yang sering digunakan}]
rsync -av --exclude='*.log' --exclude='.cache/' sumber/ tujuan/
rsync -av --exclude-from='/etc/rsync-excludes.txt' sumber/ tujuan/
rsync -av --log-file=/var/log/rsync.log sumber/ tujuan/
rsync -av --progress sumber/ tujuan/
\end{lstlisting}

\subsection*{Praktek 12.2: Sinkronisasi Direktori dengan rsync}

Tujuan praktek ini adalah memahami perilaku \cmd{rsync} dengan berbagai opsi,
termasuk efek \texttt{--delete} dan mekanisme hard link pada snapshot.

\begin{enumerate}
    \item Jalankan sinkronisasi pertama:
\begin{lstlisting}[language=bash, caption={Sinkronisasi pertama dengan rsync}]
cd ~/lab-os/chapter12-backup
rsync -avn data-sumber/ arsip-rsync/
rsync -av data-sumber/ arsip-rsync/
ls -la arsip-rsync/
\end{lstlisting}
    \textit{Amati:} apa perbedaan output dry run (\texttt{-n}) dengan eksekusi
    sebenarnya?

    \item Tambahkan file baru, hapus satu file, lalu amati efek \texttt{--delete}:
\begin{lstlisting}[language=bash, caption={Mengamati efek --delete pada rsync}]
echo "file baru penting" > data-sumber/dokumen/baru.txt
rm data-sumber/log/app.log

rsync -av --delete data-sumber/ arsip-rsync/
ls arsip-rsync/dokumen/
ls arsip-rsync/log/
\end{lstlisting}
    \textit{Amati:} apakah \file{app.log} ikut terhapus dari tujuan? Apakah
    \file{baru.txt} muncul di tujuan?

    \item Buat snapshot pertama dan kedua menggunakan \texttt{--link-dest}:
\begin{lstlisting}[language=bash, caption={Membuat dua snapshot berbasis hard link}]
mkdir -p snapshots/snap-1
rsync -av data-sumber/ snapshots/snap-1/

echo "isi baru setelah snap-1" > data-sumber/dokumen/baru.txt

mkdir -p snapshots/snap-2
rsync -av --link-dest=snapshots/snap-1 \
    data-sumber/ snapshots/snap-2/
\end{lstlisting}

    \item Verifikasi hard link dengan membandingkan inode:
\begin{lstlisting}[language=bash, caption={Verifikasi hard link antar snapshot}]
ls -lai snapshots/snap-1/dokumen/
ls -lai snapshots/snap-2/dokumen/
du -sh snapshots/snap-1/ snapshots/snap-2/
du -sh snapshots/
\end{lstlisting}
    \textit{Amati:} bandingkan nomor inode (kolom pertama) untuk file \file{laporan.txt}
    di snap-1 dan snap-2. Apakah sama? File \file{baru.txt} yang diubah
    seharusnya punya inode berbeda karena isinya berubah.

    \item Bersihkan snapshot setelah praktek:
\begin{lstlisting}[language=bash, caption={Membersihkan direktori snapshot}]
rm -rf arsip-rsync/ snapshots/
\end{lstlisting}
\end{enumerate}

\begin{challengebox}
Buat skrip Bash (mengacu ke Bab~7) bernama \file{rsync-backup.sh} di direktori
lab yang menjalankan \cmd{rsync} dari direktori \dir{data-sumber/} ke direktori
bernama sesuai tanggal hari ini dalam format \texttt{YYYY-MM-DD}. Gunakan
\cmd{tee} dari Bab~3 untuk menulis output rsync ke layar sekaligus ke file log
bernama \file{backup-TANGGAL.log}. Jadikan skrip ini executable dengan
\cmd{chmod +x} dan jalankan untuk memverifikasi hasilnya.
\end{challengebox}

%% ========================================================================
\section{tar dan dd untuk Arsip}
\label{sec:tar-dd-backup}
%% ========================================================================

\cmd{tar} (tape archiver) adalah alat pengarsipan paling klasik di Linux.
Meskipun namanya berasal dari era pita magnetik, \cmd{tar} sangat umum digunakan
untuk mengemas banyak file dan direktori menjadi satu file arsip yang bisa
dikompresi dan dipindahkan dengan mudah.

Opsi \texttt{-c} membuat arsip baru, \texttt{-f} menentukan nama file output,
dan opsi kompresi dipilih berdasarkan kebutuhan. \texttt{-z} menggunakan gzip:
cepat, kompresi sedang, ekstensi \texttt{.tar.gz}. \texttt{-j} menggunakan
bzip2: lebih lambat, kompresi lebih baik, ekstensi \texttt{.tar.bz2}. \texttt{-J}
menggunakan xz: paling lambat, kompresi terbaik, ekstensi \texttt{.tar.xz}.

\begin{lstlisting}[language=bash, caption={Membuat arsip backup dengan tar}]
tar -czf backup-$(date +%F).tar.gz data-sumber/
tar -cjf backup-$(date +%F).tar.bz2 data-sumber/
tar -cJf backup-$(date +%F).tar.xz data-sumber/

tar -czf backup.tar.gz --exclude='*.log' --exclude='.cache' data-sumber/

tar -tzf backup-$(date +%F).tar.gz
tar -tzf backup-$(date +%F).tar.gz | grep "\.conf$"

tar -xzf backup-$(date +%F).tar.gz -C /tmp/restore-test/
tar -xzf backup-$(date +%F).tar.gz -C /tmp/ data-sumber/konfigurasi/app.conf
\end{lstlisting}

Opsi \texttt{-t} menampilkan daftar isi arsip tanpa mengekstrak. Opsi \texttt{-x}
mengekstrak. Opsi \texttt{-C} menentukan direktori tujuan ekstraksi. Mengekstrak
hanya file atau direktori tertentu dilakukan dengan menyebutkan path relatif
di akhir perintah.

Konstruksi \texttt{\$(date +\%F)} menyisipkan tanggal dalam format
\texttt{YYYY-MM-DD} ke nama file secara otomatis, sehingga setiap backup memiliki
nama unik dan bisa diidentifikasi kapan dibuat.

Untuk backup incremental dengan \cmd{tar}, gunakan opsi \texttt{--listed-incremental}
yang memanfaatkan file snapshot untuk melacak perubahan.

\begin{lstlisting}[language=bash, caption={Backup incremental dengan snapshot file}]
tar -czf full-$(date +%F).tar.gz \
    --listed-incremental=snapshot.snar \
    data-sumber/

tar -czf incr-$(date +%F-%H%M).tar.gz \
    --listed-incremental=snapshot.snar \
    data-sumber/

tar -czf full-paksa-$(date +%F).tar.gz \
    --listed-incremental=/dev/null \
    data-sumber/
\end{lstlisting}

Perintah pertama menggunakan file snapshot yang belum ada, sehingga secara
otomatis membuat full backup dan menginisialisasi file snapshot. Perintah
berikutnya dengan file snapshot yang sama menghasilkan incremental backup.
Menggunakan \texttt{/dev/null} sebagai snapshot file memaksa full backup terlepas
apakah file snapshot sudah ada.

\cmd{dd} bekerja di level yang berbeda dari \cmd{tar}. Alih-alih memahami
struktur file, \cmd{dd} menyalin data byte per byte dari sumber ke tujuan.
Ini membuatnya tepat untuk backup disk atau partisi secara keseluruhan, termasuk
MBR dan struktur partisi.

\begin{lstlisting}[language=bash, caption={Penggunaan dd untuk backup disk}]
sudo dd if=/dev/sda of=/backup/sda-image.img bs=4M status=progress

sudo dd if=/dev/sda bs=4M status=progress | \
    gzip -9 > /backup/sda-$(date +%F).img.gz

sudo dd if=/backup/sda-image.img of=/dev/sdb bs=4M status=progress
\end{lstlisting}

Opsi \texttt{if} adalah input file (sumber), \texttt{of} adalah output file
(tujuan), \texttt{bs} adalah ukuran blok per operasi baca-tulis (4M memberikan
performa yang baik), dan \texttt{status=progress} menampilkan kemajuan real-time.

\begin{warningbox}
\cmd{dd} tidak meminta konfirmasi sebelum menimpa data. Menukar \texttt{if} dan
\texttt{of} secara tidak sengaja akan menghancurkan disk sumber. Selalu verifikasi
dua kali nama device dengan \cmd{lsblk} sebelum menjalankan \cmd{dd} ke disk.
\end{warningbox}

\subsection*{Praktek 12.3: Buat Arsip Backup dan Verifikasi Integritasnya}

Tujuan praktek ini adalah membuat arsip backup penuh dan incremental,
kemudian memverifikasi integritas arsip menggunakan checksum.

\begin{enumerate}
    \item Buat arsip full backup dan simpan checksumnya:
\begin{lstlisting}[language=bash, caption={Full backup dan pembuatan checksum}]
cd ~/lab-os/chapter12-backup
mkdir -p arsip-tar

tar -czf arsip-tar/full-$(date +%F).tar.gz data-sumber/
ls -lh arsip-tar/

md5sum arsip-tar/full-$(date +%F).tar.gz > arsip-tar/full-$(date +%F).md5
sha256sum arsip-tar/full-$(date +%F).tar.gz > arsip-tar/full-$(date +%F).sha256
cat arsip-tar/full-$(date +%F).md5
\end{lstlisting}
    \textit{Amati:} berapa ukuran arsip dibandingkan direktori sumber
    (\cmd{du -sh data-sumber/})?

    \item Periksa isi arsip dan coba ekstrak satu file:
\begin{lstlisting}[language=bash, caption={Memeriksa isi arsip tar}]
tar -tzf arsip-tar/full-$(date +%F).tar.gz
tar -tzf arsip-tar/full-$(date +%F).tar.gz | grep "\.conf"
mkdir -p /tmp/tar-coba
tar -xzf arsip-tar/full-$(date +%F).tar.gz \
    -C /tmp/tar-coba \
    data-sumber/konfigurasi/app.conf
cat /tmp/tar-coba/data-sumber/konfigurasi/app.conf
rm -rf /tmp/tar-coba
\end{lstlisting}

    \item Verifikasi integritas arsip menggunakan checksum:
\begin{lstlisting}[language=bash, caption={Verifikasi integritas dengan checksum}]
md5sum -c arsip-tar/full-$(date +%F).md5
sha256sum -c arsip-tar/full-$(date +%F).sha256

tar -tzf arsip-tar/full-$(date +%F).tar.gz > /dev/null && \
    echo "Arsip VALID" || echo "Arsip RUSAK"
\end{lstlisting}
    \textit{Amati:} apa arti pesan \texttt{OK} pada output \cmd{md5sum -c}?

    \item Buat arsip incremental setelah menambahkan file baru:
\begin{lstlisting}[language=bash, caption={Backup incremental dengan snapshot}]
tar -czf arsip-tar/full-snap-$(date +%F).tar.gz \
    --listed-incremental=arsip-tar/snapshot.snar \
    data-sumber/

echo "file ditambahkan setelah full backup" > data-sumber/dokumen/tambahan.txt
echo "konfigurasi baru" >> data-sumber/konfigurasi/app.conf

tar -czf arsip-tar/incr-$(date +%F-%H%M).tar.gz \
    --listed-incremental=arsip-tar/snapshot.snar \
    data-sumber/

ls -lh arsip-tar/*.tar.gz
\end{lstlisting}
    \textit{Amati:} bandingkan ukuran arsip full-snap dan arsip incremental.
    Seberapa besar penghematan dari backup incremental?
\end{enumerate}

\begin{challengebox}
Buat arsip incremental sesi kedua. Tambahkan beberapa file baru ke
\dir{data-sumber/dokumen/} dan ubah isi \file{app.conf}. Kemudian jalankan
\cmd{tar} dengan file snapshot yang sama untuk menghasilkan incremental backup
kedua. Bandingkan ukuran ketiga arsip (full, incr-1, incr-2). Kemudian verifikasi
isi setiap arsip menggunakan \cmd{tar -tzf} untuk memastikan setiap arsip
hanya mengandung file yang sesuai dengan jenisnya.
\end{challengebox}

%% ========================================================================
\section{Backup Terjadwal dengan cron}
\label{sec:scheduled-backup}
%% ========================================================================

Backup yang dilakukan secara manual tidak dapat diandalkan. Manusia lupa,
sibuk, atau sakit. \cmd{cron} adalah daemon penjadwalan Linux yang menjalankan
perintah atau skrip secara otomatis pada waktu yang telah ditentukan, seperti
alarm yang tidak pernah lupa berbunyi.

Format jadwal cron menggunakan lima field yang disebut \emph{cron expression}.
Setiap field mewakili satuan waktu yang berbeda.

\begin{lstlisting}[language=bash, caption={Format cron expression dan contoh jadwal}]
# menit jam hari-bulan bulan hari-minggu perintah
# (0-59) (0-23) (1-31) (1-12) (0-7, 0&7=Minggu)

0 2 * * *       /usr/local/bin/backup-harian.sh
0 2 * * 0       /usr/local/bin/backup-mingguan.sh
0 1 1 * *       /usr/local/bin/backup-bulanan.sh
30 22 * * 1-5   /usr/local/bin/backup-hari-kerja.sh
*/15 * * * *    /usr/local/bin/sync-kritis.sh
\end{lstlisting}

Bintang berarti ``setiap nilai yang mungkin''. \texttt{*/15} di field menit
berarti setiap 15 menit. \texttt{1-5} di field hari-minggu berarti Senin sampai
Jumat. Perintah pertama berjalan setiap hari pukul 02:00. Perintah kedua berjalan
setiap Minggu pukul 02:00.

Untuk mengedit crontab milik pengguna saat ini, gunakan \cmd{crontab -e}.
Untuk menampilkannya, gunakan \cmd{crontab -l}.

\begin{lstlisting}[language=bash, caption={Mengelola crontab}]
crontab -l
crontab -e
sudo crontab -e
ls /etc/cron.daily/ /etc/cron.weekly/ /etc/cron.monthly/
\end{lstlisting}

Direktori \dir{/etc/cron.daily/}, \dir{/etc/cron.weekly/}, dan
\dir{/etc/cron.monthly/} berisi skrip yang dijalankan otomatis pada frekuensi
yang sesuai namanya oleh \cmd{anacron}, yang menangani sistem yang tidak berjalan
24 jam terus-menerus.

Skrip backup yang dijalankan cron perlu mencatat aktivitasnya ke file log. Ini
penting karena cron berjalan tanpa terminal, sehingga tidak ada cara melihat
output langsung.

\begin{longlisting}[language=bash, caption={Contoh skrip backup dengan logging}]
#!/bin/bash

BACKUP_DIR="$HOME/lab-os/chapter12-backup/arsip-otomatis"
SOURCE="$HOME/lab-os/chapter12-backup/data-sumber"
DATE=$(date +%Y-%m-%d-%H%M)
LOG="$HOME/lab-os/chapter12-backup/backup.log"

mkdir -p "$BACKUP_DIR"

echo "[$(date)] Mulai backup $DATE" >> "$LOG"

if rsync -av --delete \
    --exclude='*.log' \
    --log-file="$LOG" \
    "$SOURCE/" \
    "$BACKUP_DIR/$DATE/"; then
    echo "[$(date)] Backup berhasil: $BACKUP_DIR/$DATE/" >> "$LOG"
else
    echo "[$(date)] GAGAL: backup $DATE" >> "$LOG"
fi

find "$BACKUP_DIR" -maxdepth 1 -type d -mtime +7 \
    -exec rm -rf {} \; 2>/dev/null
echo "[$(date)] Selesai" >> "$LOG"
\end{longlisting}

Baris \cmd{find} di akhir skrip menerapkan retensi otomatis dengan menghapus
direktori backup yang lebih lama dari 7 hari, sehingga ruang disk tidak habis
seiring waktu.

\subsection*{Praktek 12.4: Jadwalkan Skrip Backup Otomatis}

Tujuan praktek ini adalah membuat skrip backup, mengujinya secara manual,
mendaftarkan ke crontab, dan memverifikasi eksekusi otomatisnya.

\begin{enumerate}
    \item Buat skrip backup yang menggunakan rsync dengan logging:
\begin{longlisting}[language=bash, caption={Membuat skrip backup otomatis}]
cd ~/lab-os/chapter12-backup

nano backup-otomatis.sh
# Tulis isi berkas berikut
#!/bin/bash
BACKUP_BASE="$HOME/lab-os/chapter12-backup/arsip-cron"
SOURCE="$HOME/lab-os/chapter12-backup/data-sumber"
DATE=$(date +%Y-%m-%d-%H%M)
LOG="$HOME/lab-os/chapter12-backup/cron-backup.log"

mkdir -p "$BACKUP_BASE/$DATE"
echo "[$(date)] Memulai backup ke $BACKUP_BASE/$DATE" >> "$LOG"

rsync -av --delete "$SOURCE/" "$BACKUP_BASE/$DATE/" >> "$LOG" 2>&1

if [ $? -eq 0 ]; then
    UKURAN=$(du -sh "$BACKUP_BASE/$DATE/" | cut -f1)
    echo "[$(date)] OK: $DATE ($UKURAN)" >> "$LOG"
else
    echo "[$(date)] GAGAL: backup $DATE" >> "$LOG"
    rmdir "$BACKUP_BASE/$DATE" 2>/dev/null
fi

find "$BACKUP_BASE" -maxdepth 1 -type d -mmin +30 -exec rm -rf {} \; 2>/dev/null

chmod +x backup-otomatis.sh
\end{longlisting}

    \item Uji skrip secara manual sebelum mendaftarkannya ke cron:
\begin{lstlisting}[language=bash, caption={Uji eksekusi manual skrip backup}]
./backup-otomatis.sh
ls -la arsip-cron/
cat cron-backup.log
\end{lstlisting}
    \textit{Amati:} apakah log mencatat timestamp yang benar? Apakah direktori
    backup berhasil dibuat?

    \item Daftarkan skrip ke crontab untuk berjalan setiap 2 menit (untuk
        keperluan pengujian):
\begin{lstlisting}[language=bash, caption={Mendaftarkan skrip ke crontab}]
crontab -l 2>/dev/null > /tmp/cron-sementara
echo "*/2 * * * * $HOME/lab-os/chapter12-backup/backup-otomatis.sh" \
    >> /tmp/cron-sementara
crontab /tmp/cron-sementara
rm /tmp/cron-sementara
crontab -l
\end{lstlisting}

    \item Tunggu beberapa menit lalu verifikasi cron telah menjalankan backup:
\begin{lstlisting}[language=bash, caption={Verifikasi eksekusi otomatis cron}]
ls -lt arsip-cron/
cat cron-backup.log
grep CRON /var/log/syslog 2>/dev/null | tail -5
\end{lstlisting}
    \textit{Amati:} berapa direktori backup yang sudah terbuat? Apakah timestamp
    pada nama direktori sesuai dengan waktu eksekusi di log?

    \item Bersihkan crontab dan direktori:
\begin{lstlisting}[language=bash, caption={Membersihkan crontab dan hasil backup}]
crontab -l | grep -v "backup-otomatis" | crontab -
crontab -l
rm -rf arsip-cron/ arsip-tar/ cron-backup.log backup.log
\end{lstlisting}
\end{enumerate}

\begin{challengebox}
Jadwalkan skrip \file{rsync-backup.sh} dari challengebox Praktek 12.2
agar berjalan setiap hari pukul 02:00 menggunakan crontab. Tambahkan ke skrip
tersebut agar output rsync disimpan ke \file{/var/log/backup.log} menggunakan
teknik \cmd{tee} dari Bab~3 (sehingga sekaligus tercetak di terminal saat
dijalankan manual dan tersimpan ke log). Tulis cron expression yang tepat dan
jelaskan setiap field-nya.
\end{challengebox}

%% ========================================================================
\section{Pemulihan Sistem}
\label{sec:backup-recovery}
%% ========================================================================

Backup yang belum pernah diuji restore-nya tidak lebih berguna dari salinan
kosong. Proses pemulihan perlu dipraktikkan sebelum dibutuhkan dalam kondisi
darurat, karena memulihkan sistem sambil panik jauh lebih sulit daripada
memulihkannya dalam kondisi tenang.

Alur restore yang aman selalu mengekstrak ke direktori sementara terlebih dahulu,
bukan langsung ke lokasi asli. Ini melindungimu jika ada sesuatu yang tidak
beres selama proses restore: data asli yang tersisa tidak akan tertimpa sebelum
kamu yakin restore berjalan benar.

Untuk arsip \cmd{tar}, langkah pertama adalah memverifikasi integritas arsip,
kemudian mengidentifikasi file yang perlu dipulihkan.

\begin{lstlisting}[language=bash, caption={Restore dari arsip tar dengan verifikasi}]
tar -tzf /backup/full-2026-05-09.tar.gz > /dev/null && \
    echo "Arsip VALID" || echo "Arsip RUSAK"

tar -tzf /backup/full-2026-05-09.tar.gz | grep "dokumen"

mkdir -p /tmp/restore-sementara
tar -xzf /backup/full-2026-05-09.tar.gz \
    -C /tmp/restore-sementara \
    data-sumber/konfigurasi/app.conf

ls /tmp/restore-sementara/data-sumber/konfigurasi/
cp /tmp/restore-sementara/data-sumber/konfigurasi/app.conf \
   ~/target-restore/konfigurasi/

rm -rf /tmp/restore-sementara
\end{lstlisting}

Untuk snapshot \cmd{rsync}, restore bahkan lebih sederhana karena setiap snapshot
adalah direktori biasa yang bisa ditelusuri dengan \cmd{ls} tanpa perlu proses
ekstraksi.

\begin{lstlisting}[language=bash, caption={Restore dari snapshot rsync}]
ls ~/backup/snapshots/
ls ~/backup/snapshots/2026-05-09/data-sumber/

cp ~/backup/snapshots/2026-05-09/data-sumber/dokumen/laporan.txt \
   ~/target-restore/dokumen/

rsync -av ~/backup/snapshots/2026-05-09/data-sumber/konfigurasi/ \
    ~/target-restore/konfigurasi/

diff ~/backup/snapshots/2026-05-08/data-sumber/app.conf \
     ~/backup/snapshots/2026-05-09/data-sumber/app.conf
\end{lstlisting}

Setelah restore, verifikasi integritas menggunakan checksum yang disimpan saat
backup dibuat.

\begin{lstlisting}[language=bash, caption={Verifikasi integritas setelah restore}]
md5sum -c /backup/full-2026-05-09.md5
sha256sum -c /backup/full-2026-05-09.sha256

diff -r /backup/snapshots/2026-05-09/data-sumber/ ~/target-restore/
echo "Exit code: $? (0 berarti identik)"
\end{lstlisting}

\cmd{diff -r} membandingkan dua direktori secara rekursif. Exit code 0 berarti
tidak ada perbedaan: restore berhasil sempurna.

\subsection*{Praktek 12.5: Simulasi Pemulihan dari Backup}

Tujuan praktek ini adalah mensimulasikan skenario kehilangan data nyata,
lalu memulihkan dari backup yang sudah ada dan memverifikasi hasilnya.

\begin{enumerate}
    \item Siapkan data sumber dan buat backup sebagai persiapan:
\begin{lstlisting}[language=bash, caption={Siapkan data dan buat backup awal}]
cd ~/lab-os/chapter12-backup
mkdir -p data-sumber/{dokumen,konfigurasi}
echo "laporan-akhir-tahun" > data-sumber/dokumen/laporan.txt
echo "catatan-rapat" > data-sumber/dokumen/catatan.txt
echo "db_host=localhost" > data-sumber/konfigurasi/app.conf
echo "cache_ttl=300" >> data-sumber/konfigurasi/app.conf

mkdir -p arsip-pemulihan snapshot-pemulihan

tar -czf arsip-pemulihan/full-backup.tar.gz data-sumber/
md5sum arsip-pemulihan/full-backup.tar.gz > arsip-pemulihan/full-backup.md5
rsync -av data-sumber/ snapshot-pemulihan/

ls -lhR data-sumber/
\end{lstlisting}

    \item Simpan checksum file-file penting untuk verifikasi nanti:
\begin{lstlisting}[language=bash, caption={Menyimpan checksum file sumber}]
md5sum data-sumber/dokumen/* data-sumber/konfigurasi/* > checksum-asli.md5
cat checksum-asli.md5
\end{lstlisting}

    \item Simulasikan kehilangan data: hapus beberapa file secara sengaja:
\begin{lstlisting}[language=bash, caption={Simulasi kehilangan data}]
rm data-sumber/dokumen/laporan.txt
rm data-sumber/konfigurasi/app.conf
echo "file tidak sah" > data-sumber/dokumen/tidak-dikenal.txt
ls -la data-sumber/dokumen/
ls -la data-sumber/konfigurasi/
\end{lstlisting}
    \textit{Amati:} kondisi data-sumber sekarang tidak konsisten: ada file yang
    hilang dan ada file asing yang tidak dikenal.

    \item Pulihkan file yang hilang dari arsip tar ke direktori sementara:
\begin{lstlisting}[language=bash, caption={Restore dari arsip tar ke direktori sementara}]
md5sum -c arsip-pemulihan/full-backup.md5

mkdir -p /tmp/restore-lap14
tar -xzf arsip-pemulihan/full-backup.tar.gz \
    -C /tmp/restore-lap14 \
    data-sumber/dokumen/laporan.txt

ls /tmp/restore-lap14/data-sumber/dokumen/
cp /tmp/restore-lap14/data-sumber/dokumen/laporan.txt \
   data-sumber/dokumen/

rm -rf /tmp/restore-lap14
\end{lstlisting}

    \item Pulihkan file konfigurasi dari snapshot rsync:
\begin{lstlisting}[language=bash, caption={Restore dari snapshot rsync}]
ls snapshot-pemulihan/konfigurasi/
cp snapshot-pemulihan/konfigurasi/app.conf \
   data-sumber/konfigurasi/
ls -la data-sumber/konfigurasi/
\end{lstlisting}

    \item Verifikasi integritas setelah pemulihan:
\begin{lstlisting}[language=bash, caption={Verifikasi checksum setelah pemulihan}]
rm data-sumber/dokumen/tidak-dikenal.txt

md5sum -c checksum-asli.md5
echo "Exit code: $? (0 berarti semua file valid)"

diff -r data-sumber/ snapshot-pemulihan/
echo "Diff exit code: $? (0 berarti identik dengan snapshot)"
\end{lstlisting}
    \textit{Amati:} apakah semua checksum menunjukkan OK? File yang dipulihkan
    harus identik dengan versi aslinya.

    \item Bersihkan seluruh direktori lab:
\begin{lstlisting}[language=bash, caption={Membersihkan direktori praktek}]
rm -rf arsip-pemulihan/ snapshot-pemulihan/
rm -f checksum-asli.md5 rencana-backup.txt backup-otomatis.sh
\end{lstlisting}
\end{enumerate}

\begin{challengebox}
Perpanjang skenario pemulihan: hapus seluruh direktori \dir{data-sumber/konfigurasi/}
sekaligus (bukan hanya satu file). Kemudian pulihkan seluruh direktori tersebut
dari snapshot menggunakan \cmd{rsync} dengan arah transfer yang dibalik (sumber
adalah snapshot, tujuan adalah direktori yang rusak). Verifikasi hasil pemulihan
dengan membandingkan checksum seluruh isi direktori menggunakan \cmd{find} dan
\cmd{md5sum}. Catat perbedaan waktu antara restore satu file dan restore satu
direktori penuh.
\end{challengebox}

%% ========================================================================
\section{Rangkuman}
\label{sec:backup-summary}
%% ========================================================================

Dalam bab ini, kamu telah mempelajari:

\begin{itemize}[noitemsep,topsep=2pt]
    \item \textbf{Strategi backup}: tiga jenis backup (full, incremental,
        differential) dengan trade-off berbeda; aturan 3-2-1 sebagai standar
        industri; serta RPO dan RTO sebagai metrik yang menentukan frekuensi
        dan infrastruktur backup.
    \item \textbf{rsync}: sinkronisasi efisien berbasis delta-transfer, opsi
        \texttt{--delete} untuk cermin direktori, dan \texttt{--link-dest} untuk
        snapshot backup berbasis hard link yang hemat ruang disk.
    \item \textbf{tar}: pengarsipan dengan pilihan kompresi gzip, bzip2, dan xz;
        backup incremental menggunakan opsi \texttt{--listed-incremental};
        serta restore file individual dari arsip.
    \item \textbf{cron}: penjadwalan otomatis dengan format cron expression lima
        field, pengelolaan crontab, dan teknik logging output skrip ke file.
    \item \textbf{Pemulihan dan verifikasi}: prosedur restore yang aman melalui
        direktori sementara, verifikasi integritas dengan \cmd{md5sum} dan
        \cmd{sha256sum}, serta pentingnya menguji restore secara berkala.
\end{itemize}

%% ========================================================================
\section{Latihan}
\label{sec:latihan-backup}
%% ========================================================================

\textbf{Instruksi Umum}: Kerjakan seluruh latihan secara mandiri. Catat langkah penting, simpan tangkapan layar bila diperlukan, lalu rangkum hasilnya sebagai dokumentasi pribadi.

\subsubsection*{Latihan 12.1 Implementasi Sistem Backup Lengkap}
Rancang dan implementasikan sistem backup untuk direktori simulasi.
\begin{enumerate}
    \item Buat struktur direktori simulasi dengan minimal 10 file yang tersebar
        di tiga subdirektori: \dir{dokumen/}, \dir{konfigurasi/}, dan \dir{media/}.
    \item Buat skrip \file{backup-harian.sh} yang menggunakan \cmd{rsync} dengan
        \texttt{--link-dest} untuk membuat snapshot harian. Skrip harus mencatat
        log dengan timestamp ke \file{backup-harian.log}.
    \item Buat skrip \file{backup-mingguan.sh} yang menggunakan \cmd{tar -czf}
        untuk membuat arsip terkompresi. Skrip harus membuat checksum MD5 dari
        setiap arsip yang dibuat.
    \item Daftarkan keduanya ke crontab dengan jadwal yang berbeda. Jalankan
        masing-masing secara manual dan verifikasi log dan output yang dihasilkan.
\item Simulasikan kehilangan file dan lakukan restore dari kedua jenis backup.
        Dokumentasikan langkah-langkah restore dan waktu yang dibutuhkan.
\end{enumerate}

\subsubsection*{Latihan 12.2 Analisis Kompresi dan Performa Backup}
Analisis trade-off antara kecepatan dan rasio kompresi.
\begin{enumerate}
    \item Buat direktori dengan tiga jenis file: teks biasa (10 file \texttt{.txt}
        masing-masing 100 baris), file konfigurasi \texttt{.conf}, dan file biner
        simulasi menggunakan \cmd{dd if=/dev/urandom of=biner.bin bs=1M count=5}.
    \item Buat tiga arsip dari direktori yang sama menggunakan gzip (\texttt{-z}),
        bzip2 (\texttt{-j}), dan xz (\texttt{-J}). Ukur waktu setiap proses
        menggunakan \cmd{time}.
    \item Bandingkan ukuran ketiga arsip dengan \cmd{ls -lh} dan hitung rasio
        kompresi masing-masing terhadap ukuran asli.
    \item Buat tabel di file \file{analisis-kompresi.txt} yang merangkum: jenis
        kompresi, waktu kompres, ukuran hasil, dan rasio kompresi.
\item Berdasarkan data tersebut, rekomendasikan kompresi yang paling tepat
        untuk: backup harian otomatis, arsip jangka panjang, dan backup file biner.
        Berikan alasan untuk setiap rekomendasi.
\end{enumerate}

\subsubsection*{Latihan 12.3 Disaster Recovery Drill}
Lakukan simulasi pemulihan bencana secara menyeluruh.
\begin{enumerate}
    \item Buat direktori ``produksi'' dengan struktur lengkap: file konfigurasi,
        dokumen, dan skrip. Buat backup penuh menggunakan \cmd{tar} dan simpan
        checksumnya.
    \item Dokumentasikan kondisi awal: daftar file, ukuran, dan checksum semua
        file menggunakan \cmd{find} dan \cmd{md5sum}.
    \item Simulasikan bencana: hapus seluruh direktori produksi dengan \cmd{rm -rf}.
    \item Catat waktu mulai pemulihan, lakukan restore lengkap dari backup, dan
        catat waktu selesai. Hitung RTO aktual.
    \item Verifikasi semua file pulih dengan benar menggunakan checksum yang
        disimpan di langkah 2.
\item Bandingkan RTO aktual dengan target yang kamu tentukan. Jika lebih lama,
        identifikasi bottleneck dan usulkan cara mempercepatnya.
\end{enumerate}

%% ========================================================================
\section{Referensi}
\label{sec:referensi-backup}
%% ========================================================================

\begin{itemize}
    \item Tridgell, Andrew dan Mackerras, Paul. \textit{The rsync Algorithm}.
        Australian National University, 1996.
        \url{https://rsync.samba.org/tech_report/}
    \item GNU Project. \textit{GNU tar Manual}.
        \url{https://www.gnu.org/software/tar/manual/}
    \item Canonical Ltd. \textit{Ubuntu Server Guide: Backups}.
        \url{https://ubuntu.com/server/docs/backups-introduction}
    \item Shotts, William E. \textit{The Linux Command Line}, 5th ed.
        No Starch Press, 2019. \url{https://linuxcommand.org/tlcl.php}
    \item Turnbull, James. \textit{The Linux Command Line for Beginners}.
        Apress, 2012.
\end{itemize}
